\documentclass{article}

% ------------------------------------------------------------------------------------
% A ZOO Deep Learning Survey -- NeurIPS 2025 format, camera-ready (non-anonymous).
%
% The official neurips_2025.sty defines the footer notice only when a track option is
% given, so the non-anonymous "final" option is combined with "main" exactly as the
% template prescribes for camera-ready copies:  \usepackage[main, final]{neurips_2025}.
% ------------------------------------------------------------------------------------
\usepackage[main,final]{neurips_2025}

\usepackage[utf8]{inputenc}
\usepackage[T1]{fontenc}
\usepackage{hyperref}
\usepackage{url}
\usepackage{booktabs}
\usepackage{amsfonts}
\usepackage{amsmath}
\usepackage{amssymb}
\usepackage{nicefrac}
\usepackage{microtype}
\usepackage{xcolor}
\usepackage{graphicx}
\usepackage{tikz}
\usepackage{pgfplots}
\pgfplotsset{compat=1.18}
\usepgfplotslibrary{groupplots}

\hypersetup{colorlinks=true, linkcolor=blue!60!black, citecolor=blue!60!black, urlcolor=blue!60!black}

% ------------------------------------------------------------------------------------
% Shared colour / marker cycle for the 8 optimizers.  Order is fixed everywhere:
%   bp  spsa  rs  np  anp  inp  danp  ff
% bp is black dashed; the seven others use an Okabe-Ito style colour-blind-safe set.
% ------------------------------------------------------------------------------------
\definecolor{zoospsa}{HTML}{E69F00}
\definecolor{zoors}{HTML}{0072B2}
\definecolor{zoonp}{HTML}{009E73}
\definecolor{zooanp}{HTML}{D55E00}
\definecolor{zooinp}{HTML}{CC79A7}
\definecolor{zoodanp}{HTML}{56B4E9}
\definecolor{zooff}{HTML}{6A3D9A}
\pgfplotscreateplotcyclelist{zoocycle}{%
  color=black,   dashed, thick, mark=none\\
  color=zoospsa, solid, mark=o,         mark size=1.3pt\\
  color=zoors,   solid, mark=square*,   mark size=1.2pt\\
  color=zoonp,   solid, mark=triangle*, mark size=1.5pt\\
  color=zooanp,  solid, mark=diamond*,  mark size=1.5pt\\
  color=zooinp,  solid, mark=pentagon*, mark size=1.5pt\\
  color=zoodanp, solid, mark=x,         mark size=1.8pt\\
  color=zooff,   solid, mark=+,         mark size=1.8pt\\
}

% Adds the 8 optimizer curves of one CSV (#2) with x column #1, in zoocycle order.
\newcommand{\zooplots}[2]{%
  \addplot table [x=#1, y=bp,   col sep=comma] {#2};
  \addplot table [x=#1, y=spsa, col sep=comma] {#2};
  \addplot table [x=#1, y=rs,   col sep=comma] {#2};
  \addplot table [x=#1, y=np,   col sep=comma] {#2};
  \addplot table [x=#1, y=anp,  col sep=comma] {#2};
  \addplot table [x=#1, y=inp,  col sep=comma] {#2};
  \addplot table [x=#1, y=danp, col sep=comma] {#2};
  \addplot table [x=#1, y=ff,   col sep=comma] {#2};
}

% Short names used in the text.
\newcommand{\method}[1]{\textsc{#1}}
\newcommand{\BP}{\method{bp}}
\newcommand{\SPSA}{\method{spsa}}
\newcommand{\RS}{\method{rs}}
\newcommand{\NP}{\method{np}}
\newcommand{\ANP}{\method{anp}}
\newcommand{\INP}{\method{inp}}
\newcommand{\DANP}{\method{danp}}
\newcommand{\FF}{\method{ff}}
\newcommand{\ghat}{\hat{\mathbf{g}}}
\newcommand{\btheta}{\boldsymbol{\theta}}
\newcommand{\bxi}{\boldsymbol{\xi}}
\newcommand{\bz}{\mathbf{z}}
\newcommand{\bv}{\mathbf{v}}
\newcommand{\bh}{\mathbf{h}}
\newcommand{\ba}{\mathbf{a}}
\newcommand{\bb}{\mathbf{b}}
\newcommand{\bW}{\mathbf{W}}
\newcommand{\bR}{\mathbf{R}}
\newcommand{\bC}{\mathbf{C}}
\newcommand{\bI}{\mathbf{I}}
\newcommand{\dL}{\delta\mathcal{L}}
\newcommand{\loss}{\mathcal{L}}
\DeclareMathOperator*{\argmin}{arg\,min}

\title{A ZOO Deep Learning Survey}

\author{%
  Francois Chaubard \\
  Stanford University \\
  \texttt{fchaubar@stanford.edu}
}

\begin{document}

\maketitle

\begin{abstract}
Zero-order optimization (ZOO) trains a network from loss evaluations alone, without
back-propagating errors. It is the natural training rule for neuromorphic and analog
hardware, for systems that contain non-differentiable components, for biologically
plausible learning, and for memory-constrained fine-tuning. Many ZOO rules have been
proposed, but they are rarely compared under one budget on the same networks. We
survey eight training rules---backpropagation as the reference, simultaneous
perturbation (SPSA), random search, four node-perturbation variants (classic,
activity-based, iterative and decorrelated), and the forward-forward algorithm---on
four standard $\approx$40k-parameter architectures: an MLP on MNIST, a convolutional
network on CIFAR-10, and a character-level LSTM and Transformer on Penn Treebank. Every
method is given the same 16-forward-pass step budget and the same one hour of
wall-clock training on a single GPU, feeds its gradient estimate to Adam, and is swept
over three probe scales, yielding 96 runs. We report final test loss, accuracy and
bits-per-character, training curves against wall-clock, sensitivity to the probe
scale, and throughput.
Iterative node perturbation (INP) is the most consistent zero-order rule: it is the
best probe method on both sequence tasks and within a point of the best on both image
tasks, where activity-based (ANP) and decorrelated (DANP) node perturbation win. At
their best evaluation the node-perturbation rules reach 96.7\% on MNIST and 62.9\% on
CIFAR-10 against 97.5\% and 69.7\% for backpropagation, but on the character-level
LSTM and Transformer the best zero-order rule still trails backpropagation by about
one bit per character (2.82 versus 1.78 and 3.10 versus 1.82 bits/char). Weight-space
perturbation (SPSA, random search) is competitive only on the sequence models, where
the number of perturbed units approaches the number of weights, and the
forward-forward readout lags far behind its own label-accumulation inference (19\%
versus 96\% on MNIST).
Code, configurations and raw logs are released at
\url{https://github.com/Fchaubard/zero-order-ablations}.
\end{abstract}

% ====================================================================================
\section{Introduction}
\label{sec:intro}

Backpropagation \citep{rumelhart1986} is the default way to train neural networks, but
it is not always available. Neuromorphic and analog accelerators implement a forward
pass natively and may have no mechanism to transport exact error signals backwards
through the substrate; systems with non-differentiable components (simulators,
hardware in the loop, discrete decisions, black-box APIs) expose only a loss value;
neuroscience asks which learning rules a brain could plausibly implement with local
information; and fine-tuning of very large models is increasingly limited by the
activation memory that backpropagation must store \citep{malladi2023}. In all of these
settings the learner can evaluate the loss at perturbed parameters or perturbed
activities, but cannot differentiate it. \emph{Zero-order optimization} (ZOO) is the
family of methods that estimate a descent direction from such evaluations alone.

The ZOO literature is old and broad: random search \citep{rastrigin1963,matyas1965,
schumer1968}, simultaneous perturbation stochastic approximation
\citep{spall1992,spall1998}, Gaussian smoothing and its convergence theory
\citep{nesterov2017,duchi2015,flaxman2005}, evolution strategies
\citep{wierstra2014,hansen2003,salimans2017,mania2018}, node perturbation
\citep{widrow1990,werfel2003,fiete2006,hiratani2022} and its recent activity-based and
decorrelated variants \citep{dalm2023}, and local learning rules such as
forward-forward \citep{hinton2022}. Each paper demonstrates its method on its own
architecture, with its own budget of forward passes, its own choice of optimizer and
its own hyper-parameter search. A practitioner who must train without backpropagation
is left with the question this paper tries to answer: \emph{which ZOO method wins, on
standard small architectures, when every method is given exactly the same budget?}

We answer it empirically. Our protocol fixes four $\approx$40k-parameter models on four
standard tasks (an MLP on MNIST, a ConvNet on CIFAR-10, and a character-level LSTM and
Transformer on Penn Treebank), gives every probe-based method 16 forward passes per
step, runs every method for one hour of wall-clock training on one GPU, feeds every
method's gradient estimate to the same Adam optimizer, and sweeps a single scalar
$\varepsilon$ that is simultaneously the learning rate and the probe scale. Model
selection uses validation loss only; test metrics are reported once.

\paragraph{Contributions.}
\begin{itemize}
  \item A single, fully specified protocol (Section~\ref{sec:methods} and
        Section~\ref{sec:setup}) under which eight training rules---backpropagation,
        SPSA, random search, classic / activity-based / iterative / decorrelated node
        perturbation and forward-forward---are compared on four architectures with
        identical forward-pass, wall-clock and hyper-parameter budgets.
  \item A faithful implementation of each published estimator on top of one
        \emph{perturbable layer} abstraction (Appendix~\ref{app:impl}) so that every
        method drives the same models, including an LSTM and a pre-LN Transformer, not
        only MLPs.
  \item The resulting 96-run grid (Table~\ref{tab:main}, Appendix~\ref{app:grid}),
        training curves against wall-clock (Figure~\ref{fig:curves}), the sensitivity
        of every method to the probe scale (Figure~\ref{fig:eps}) and the throughput of
        every method (Table~\ref{tab:throughput}), with code and logs released.
\end{itemize}

Table~\ref{tab:main} previews the main result: iterative node perturbation is the most
consistent zero-order rule (best on both sequence tasks, mean rank 2.0 of 7 across
the four tasks), activity-based and decorrelated node perturbation win the two image
tasks, and the best zero-order rule matches backpropagation's accuracy on MNIST to
within a point but trails it by about one bit per character on the character-level
LSTM and Transformer after one hour.

% ====================================================================================
\section{Related Work}
\label{sec:related}

\paragraph{SPSA, Gaussian smoothing and evolution strategies.}
Simultaneous perturbation stochastic approximation \citep{spall1992,spall1998}
estimates a gradient from two loss evaluations at $\btheta\pm\varepsilon\bz$ with a
random direction $\bz$; with Gaussian $\bz$ the estimator is the gradient of the
Gaussian-smoothed loss \citep{nesterov2017}, and its complexity in the convex
\citep{duchi2015} and bandit \citep{flaxman2005} settings is well understood (see
\citet{liu2020primer} for a survey). Evolution strategies use the same estimator with
many directions, fitness shaping and natural-gradient corrections
\citep{wierstra2014,hansen2003,hansen2016}; \citet{salimans2017} and \citet{mania2018}
showed that they are competitive with policy-gradient reinforcement learning, which is
itself a zero-order method with respect to the environment \citep{williams1992}.
Adaptive-moment updates on top of zero-order estimates were studied by
\citet{chen2019zoadamm}. Our \SPSA{} and \RS{} baselines are the two-point and
best-of-$K$ instances of this family.

\paragraph{Node perturbation.}
Rather than perturbing weights, node perturbation injects noise into the units of the
network and correlates the resulting change of the loss with the noise, which yields a
weight update that is local to each layer
\citep{widrow1990,werfel2003,werfel2005,fiete2006}. \citet{hiratani2022} analysed its
stability and scalability and showed that it suffers from weight diffusion in
mismatched networks. \citet{dalm2023} reformulated node perturbation as a directional
derivative (\emph{iterative} NP), proposed the cheaper \emph{activity-based} NP that
uses the measured activity change rather than the injected noise, and combined it with
the input decorrelation rule of \citet{ahmad2023constrained}, which was also shown to
accelerate backpropagation itself \citep{dalm2024,ahmad2024}. We implement all four
variants.

\paragraph{Forward gradients.}
Forward-mode automatic differentiation gives an unbiased directional derivative in a
single forward pass; \citet{baydin2022} proposed training with a random direction, and
\citet{ren2023} scaled it with local losses. Forward gradients require a
Jacobian-vector product through the network and therefore do not apply to the
black-box and hardware settings that motivate this survey; we do not include them.

\paragraph{Decorrelation.}
Correlated layer inputs make the loss landscape ill-conditioned; decorrelating inputs
with a local rule \citep{ahmad2023constrained} removes this and was reported to make
node perturbation competitive with backpropagation \citep{dalm2023} and to speed up
backpropagation \citep{dalm2024,ahmad2024}. Our \DANP{} baseline uses exactly this
rule.

\paragraph{Forward-forward and other local rules.}
The forward-forward algorithm \citep{hinton2022} replaces the backward pass by two
forward passes with positive and negative data and trains each layer to separate
them with a local ``goodness'' objective. Other backprop-free rules, such as feedback
alignment \citep{lillicrap2016}, direct feedback alignment \citep{nokland2016} and
synthetic gradients \citep{jaderberg2017}, still require a backward signal that is
learned or random; they are outside our zero-order scope and we include only
forward-forward as the representative local rule.

\paragraph{Zero-order fine-tuning of large language models.}
MeZO \citep{malladi2023} showed that two-point SPSA with in-place perturbation can
fine-tune billion-parameter language models with inference-level memory; DeepZero
\citep{chen2024deepzero} scaled zero-order training from scratch with sparsity and
coordinate-wise estimates; and \citet{zhang2024revisiting} benchmarked the resulting
estimators and optimizers for memory-efficient fine-tuning. Those works vary the
estimator on one large model; we vary the estimator \emph{and} the architecture at a
scale where every method can be run to its one-hour limit.

% ====================================================================================
\section{Methods}
\label{sec:methods}

\subsection{Problem setup and the shared update rule}
\label{sec:setup-rule}

Let $\btheta\in\mathbb{R}^{P}$ be the parameters of a network with $L$ perturbable
layers, and let $\loss(\btheta)=\frac{1}{N}\sum_{n=1}^{N}\ell_n(\btheta)$ be the
mini-batch cross-entropy loss, where $\ell_n$ is the per-example loss (for language
modelling, the per-sequence mean of the 64 next-character cross-entropies). Every method
in this paper produces an estimate $\ghat\approx\nabla_{\btheta}\loss$ of the gradient
and writes it into the gradient slot of every parameter; backpropagation writes the
exact gradient, the zero-order methods write their estimator, and forward-forward
writes the gradient of its local objectives. The estimate is then consumed by the same
optimizer, Adam \citep{kingma2015} with $(\beta_1,\beta_2)=(0.9,0.99)$ and no schedule:
\begin{equation}
  \mathbf{m}\leftarrow\beta_1\mathbf{m}+(1-\beta_1)\,\ghat,\qquad
  \mathbf{v}\leftarrow\beta_2\mathbf{v}+(1-\beta_2)\,\ghat^{2},\qquad
  \btheta\leftarrow\btheta-\varepsilon\,\frac{\hat{\mathbf{m}}}{\sqrt{\hat{\mathbf{v}}}+10^{-8}},
  \label{eq:adam}
\end{equation}
with the usual bias corrections $\hat{\mathbf{m}},\hat{\mathbf{v}}$. The learning rate
$\varepsilon$ is the only hyper-parameter, and it \emph{doubles as the probe scale}: it
is the standard deviation of the parameter perturbation for \SPSA{} and \RS{} and the
standard deviation $\sigma$ of the node noise for the node-perturbation family. All
probes are drawn from $\mathcal{N}(\mathbf{0},\bI)$ and scaled by $\varepsilon$. Tying
the two removes a two-dimensional grid search per method and expresses the natural
constraint that a probe cannot be finer than the step it informs.

Because Adam normalises $\ghat$ by a running estimate of its own magnitude,
Eq.~\eqref{eq:adam} is approximately invariant to a constant rescaling of the
estimator. The normalisation constants of the estimators below therefore influence the
first few hundred steps and the relative weighting of estimates across steps, but not
the asymptotic step size; this is deliberate and makes the comparison insensitive to
the different conventions in the literature.

\subsection{Optimizers}
\label{sec:optimizers}

\paragraph{SPSA (central difference, Gaussian probes).}
With $K=8$ antithetic pairs of probes $\bz_k\sim\mathcal{N}(\mathbf{0},\bI_P)$ the
estimator is the two-point central difference
\begin{equation}
  \ghat_{\SPSA}
  = \frac{1}{K}\sum_{k=1}^{K}\frac{\loss(\btheta+\varepsilon\bz_k)-\loss(\btheta-\varepsilon\bz_k)}{2\varepsilon}\,\bz_k ,
  \qquad K=8,
  \label{eq:spsa}
\end{equation}
which is the Gaussian-smoothing estimator of \citet{nesterov2017} and the
multi-direction form of SPSA \citep{spall1992} with Gaussian rather than Rademacher
directions. Each pair costs two forward passes on the batch-mean loss, so the 16-pass
budget buys $K=8$ directions. The implementation mirrors the in-place perturbation
of MeZO \citep{malladi2023}: the parameters are restored exactly from a flat copy
after each pair. No clean pass is spent; the training loss reported for \SPSA{} is
the mean of the 16 probe losses.

\paragraph{Random search (jumping).}
Random search \citep{rastrigin1963,matyas1965,schumer1968} evaluates the clean loss
$\loss_0=\loss(\btheta)$ and $K=15$ candidates $\btheta+\varepsilon\bz_k$,
$\bz_k\sim\mathcal{N}(\mathbf{0},\bI_P)$, and keeps the best one,
$k^{*}=\argmin_k\loss(\btheta+\varepsilon\bz_k)$. Instead of jumping to the candidate
directly, it is expressed in the shared update rule through the pseudo-gradient
\begin{equation}
  \ghat_{\RS}=
  \begin{cases}
    \dfrac{\loss_{k^{*}}-\loss_0}{\varepsilon}\,\bz_{k^{*}} & \text{if }\loss_{k^{*}}<\loss_0,\\[1.2ex]
    \mathbf{0} & \text{otherwise},
  \end{cases}
  \label{eq:rs}
\end{equation}
which points along the improving direction (the difference is negative) and is zero
when no candidate improves on the clean loss; Adam then takes a step of magnitude
$\approx\varepsilon$ along $\bz_{k^{*}}$. The acceptance rate (fraction of steps with a
non-zero update) is logged. Adam's moment estimates are updated on rejected steps as
well (with a zero gradient), so a rejected step still moves the parameters along the
decaying momentum.

\paragraph{Node perturbation (NP).}
Classic node perturbation \citep{widrow1990,werfel2003,fiete2006,hiratani2022} perturbs
the units rather than the weights. For layer $l$ with input $\bh_{l-1}$,
pre-activation $\ba_l=\bW_l\bh_{l-1}+\bb_l$ and post-activation
$\bh_l=f(\ba_l)$, a noisy forward pass injects Gaussian noise into \emph{every} layer
and propagates it:
\begin{equation}
  \tilde{\ba}_l=\bW_l\tilde{\bh}_{l-1}+\bb_l+\sigma\,\bxi_l,\qquad
  \bxi_l\sim\mathcal{N}(\mathbf{0},\bI),\qquad \tilde{\bh}_l=f(\tilde{\ba}_l),
  \label{eq:np-pass}
\end{equation}
with $\sigma=\varepsilon$. Writing $\dL_n=\tilde\ell_n-\ell_n$ for the change of the
per-example loss, the estimator for the weights and biases of layer $l$ is
\begin{equation}
  \ghat_{\bW_l}^{\NP}=\frac{1}{N}\sum_{n=1}^{N}\frac{\dL_n}{\sigma}\,\bxi_{l,n}\,\bh_{l-1,n}^{\top},
  \qquad
  \ghat_{\bb_l}^{\NP}=\frac{1}{N}\sum_{n=1}^{N}\frac{\dL_n}{\sigma}\,\bxi_{l,n},
  \label{eq:np}
\end{equation}
i.e.\ the outer product of the noise-weighted loss change with the \emph{clean} layer
input (for convolutions, the matching correlation over unfolded patches; for
embeddings, the one-hot input). With one clean pass and 15 noisy passes the estimator
is averaged over 15 independent noise draws. Note that $\dL_n/\sigma\cdot\bxi_l$ equals
$\dL_n\,\varepsilon_l/\sigma^{2}$ with $\varepsilon_l=\sigma\bxi_l$, the form used by
\citet{dalm2023}; the implementation uses exactly Eq.~\eqref{eq:np} with the batch
average over the per-example loss changes, so that $\ghat$ estimates the gradient of
the batch-mean loss on the same scale as backpropagation. For a layer applied several
times per forward pass (the LSTM gate layer, once per time step) the same $\dL_n$
weights every call.

\paragraph{Activity-based node perturbation (ANP).}
Because the noise injected in early layers propagates, the actual perturbation of
layer $l$ in Eq.~\eqref{eq:np-pass} is not $\sigma\bxi_l$ but the measured activity
difference $\delta\ba_l=\tilde{\ba}_l-\ba_l$. Activity-based NP \citep{dalm2023}
replaces the injected noise by this measured difference and normalises by the norm of
the total activity perturbation across all layers:
\begin{equation}
  \ghat_{\bW_l}^{\ANP}=\frac{1}{N}\sum_{n=1}^{N}
    c\,\frac{\dL_n}{\|\delta\ba_n\|^{2}}\;\delta\ba_{l,n}\,\bh_{l-1,n}^{\top},
  \qquad
  \|\delta\ba_n\|^{2}=\sum_{l'=1}^{L}\|\delta\ba_{l',n}\|^{2},
  \label{eq:anp}
\end{equation}
where $c$ is a normalisation constant ($c=\sum_l N_l$, the total number of perturbed
units, in \citet{dalm2023}), with the bias analogue obtained by dropping
$\bh_{l-1,n}^{\top}$. For a single perturbed unit the estimator reduces to
Eq.~\eqref{eq:np}; in general it corrects for the propagated perturbation that
classic NP ignores. Following the authors' reference implementation, $c=\sum_l N_l$
counts every perturbed pre-activation unit of the network (all layers, all time steps
or positions; the input is not counted) and the norm is taken per example over the
same set of units.

\paragraph{Iterative node perturbation (INP).}
Iterative NP \citep{dalm2023} removes the propagation problem differently: each noisy
pass perturbs \emph{one} layer only, so that the loss change $\dL^{(l)}_n$ is
attributable to that layer and the estimator is a directional derivative along the
injected noise $\bv_{l,n}=\sigma\bxi_{l,n}$,
\begin{equation}
  \ghat_{\bW_l}^{\INP}=\frac{1}{N}\sum_{n=1}^{N}
    N_l\,\frac{\dL^{(l)}_n}{\|\bv_{l,n}\|^{2}}\;\bv_{l,n}\,\bh_{l-1,n}^{\top},
  \label{eq:inp}
\end{equation}
with $N_l$ the number of perturbed units of layer $l$ (Eq.~3 of \citet{dalm2023}),
at the price of one pass per layer: with $L$ perturbable layers, $L$ passes are needed
for one estimate of the full gradient. The 15 noisy passes of the budget give
$\lfloor 15/L\rfloor$ draws per layer, i.e.\ 5 for the MLP and the LSTM ($L=3$), 3
for the ConvNet ($L=5$) and exactly 1 for the Transformer ($L=15$: two embeddings,
twelve projection layers of the two blocks and the readout), so every architecture uses
16 passes (Table~\ref{tab:budget}).

\paragraph{Decorrelated ANP (DANP).}
Decorrelated ANP \citep{dalm2023} is ANP in which the input of every layer is first
decorrelated by a square matrix $\bR_l$ that is \emph{not} a parameter of the model:
$\hat{\bh}_{l-1}=\bR_l\bh_{l-1}$, $\ba_l=\bW_l\hat{\bh}_{l-1}+\bb_l$, with
$\bR_l$ initialised to the identity and trained by the local rule of
\citet{ahmad2023constrained},
\begin{equation}
  \bR_l\leftarrow\bR_l-\eta_R\,\bigl(\hat{\bC}_l-\operatorname{diag}\hat{\bC}_l\bigr)\,\bR_l,
  \qquad
  \hat{\bC}_l=\frac{1}{M}\sum_{m=1}^{M}\hat{\bh}_{l-1,m}\hat{\bh}_{l-1,m}^{\top},
  \label{eq:decor}
\end{equation}
where $\hat{\bC}_l$ is the batch covariance of the decorrelated inputs recorded on the
clean pass ($M$ rows: examples, or examples $\times$ positions for sequence and
convolutional layers, whose decorrelated vector is the unfolded $C_{\text{in}}k^{2}$
patch). The ANP estimator \eqref{eq:anp} is then applied with $\hat{\bh}_{l-1}$ in
place of $\bh_{l-1}$. Embedding layers, whose inputs are one-hot, are not
decorrelated. The decorrelation rate is fixed at $\eta_R=10^{-3}$ (the value of
\citet{dalm2023}, Appendix~D, and the default of their code) and is not tied to
$\varepsilon$; $\bR_l$ is updated once per step, after the Adam step, from the inputs
recorded on that step's clean pass. The $\bR_l$ add no capacity (they compose linearly
with $\bW_l$) and are not counted in the parameter budget.

\paragraph{Forward-forward (FF).}
The forward-forward algorithm \citep{hinton2022} trains each layer with a local
objective and never propagates an error signal between layers. For a layer with
activity $\bh_l$ the goodness is $G=\operatorname{mean}_j h_{l,j}^{2}$, and the local
loss on a positive/negative pair is
\begin{equation}
  \ell_{\FF}=\log\bigl(1+e^{\,\theta-G_{\text{pos}}}\bigr)+\log\bigl(1+e^{\,G_{\text{neg}}-\theta}\bigr),
  \qquad \theta=2,
  \label{eq:ff}
\end{equation}
whose gradient with respect to the layer's own parameters is written into the
gradient slot (local backward within the layer only). Activities are
$\ell_2$-normalised before being passed to the next layer so that only their
direction, not their goodness, is forwarded. For classification, positive data are
inputs with the correct label overlaid on the first 10 pixels and negative data carry
a wrong label; a softmax readout, trained with the same optimizer on the normalised
top-layer activity of a pass with a neutral (uniform) label, provides the test loss
and accuracy. For character-level language modelling, positive sequences are the real
text, negative sequences have 30\,\% of their characters replaced at random, and the
readout predicts the next character at every position from the normalised top-layer
activity. The positive, negative and neutral-label inputs are processed together as one
batch of $3N$ examples, so a step costs one positive and one negative pass over the
blocks; the readout is trained on detached features and therefore never shapes the
hidden layers. A single Adam instance covers all parameters: every parameter receives
exactly one local gradient per step. Because the LSTM output $\sigma(o)\tanh(c)$ lies
in $(-1,1)$, its goodness cannot exceed 1 and the threshold $\theta=2$ would saturate
every unit for positive and negative data alike; for that block we use $\theta=0.25$,
a quarter of the attainable range. We report both the readout accuracy and Hinton's
label-accumulation accuracy (the label whose overlay yields the highest total
goodness).

\paragraph{Backpropagation (reference).}
One forward and one backward pass produce the exact gradient
$\ghat_{\BP}=\nabla_{\btheta}\loss$, which is fed to the same Adam update with the same
$\varepsilon$. It is the upper reference, not a zero-order method.

\subsection{Forward-pass budget}
\label{sec:budget}

Every probe-based method is given exactly 16 forward passes per step. The passes are
spent differently (Table~\ref{tab:budget}): \SPSA{} uses all 16 for 8 antithetic pairs
and has no clean pass; \RS{} and the node-perturbation family use one clean pass to
record the clean activities (and, for \RS{}, the clean loss) plus 15 probe passes.
Forward-forward uses two passes (positive and negative) and backpropagation one forward
plus one backward pass; they are not probe methods and are included as references for
the two ends of the spectrum. Because the number of passes per step is fixed, the
number of steps per hour differs between methods (Section~\ref{sec:throughput}).

\begin{table}[t]
  \caption{How each method spends its forward-pass budget per step. $L$ is the number
  of perturbable layers (3 for the MLP and the LSTM, 5 for the ConvNet, 15 for the
  Transformer).}
  \label{tab:budget}
  \centering
  \footnotesize
  \setlength{\tabcolsep}{4pt}
  \begin{tabular}{lcccl}
    \toprule
    Method & Clean passes & Probe passes & Backward & Estimates per step \\
    \midrule
    \BP{}   & 1 & 0  & 1 full      & exact gradient \\
    \SPSA{} & 0 & 16 & --          & 8 central differences, averaged \\
    \RS{}   & 1 & 15 & --          & best of 15 candidates \\
    \NP{}   & 1 & 15 (all layers) & --          & 15 noise draws, averaged \\
    \ANP{}  & 1 & 15 (all layers) & --          & 15 noise draws, averaged \\
    \INP{}  & 1 & 15 (one layer each) & --      & $15/L$ draws per layer \\
    \DANP{} & 1 & 15 (all layers) & --          & 15 noise draws, averaged; local $\bR_l$ update \\
    \FF{}   & 2 (pos.\ + neg.) & 0 & per layer, local & local goodness gradients \\
    \bottomrule
  \end{tabular}
\end{table}

% ====================================================================================
\section{Experimental Setup}
\label{sec:setup}

\paragraph{Tasks and architectures.}
Table~\ref{tab:tasks} lists the four tasks. All models have $\approx$40k parameters so
that a run fits comfortably within the budget and the methods are compared at the same
scale. Weights are Xavier-uniform initialised \citep{glorot2010} with zero biases (the
LSTM forget-gate bias is set to 1); layer normalisation \citep{ba2016} in the
Transformer carries no affine parameters, so that every parameter of every model
lives in a perturbable layer. MNIST \citep{lecun1998} and CIFAR-10 \citep{krizhevsky2009}
use the last 5k images of the official training sets as validation; inputs are
standardised (per channel for CIFAR-10). Penn Treebank \citep{marcus1993} is used at
the character level \citep{mikolov2012} with the standard train/valid/test files and
a vocabulary built from the training text; a word-level embedding would by itself
exceed the 40k-parameter budget. Classification uses batches of 128 images; language
modelling uses batches of 64 sequences of 64 characters. There is no augmentation,
dropout or weight decay.

\begin{table}[t]
  \caption{Tasks, data splits (train/validation/test) and architectures. All
  hidden nonlinearities are ReLU; convolutions are $3{\times}3$ with $2{\times}2$
  max-pooling. The LSTM \citep{hochreiter1997} has one layer of 80 units; the
  Transformer \citep{vaswani2017} uses two pre-LN blocks, model width 48, 4 heads,
  feed-forward width 64, learned positional embeddings and context length $T=64$.
  $V$ is the character vocabulary size. \emph{Units} is the number of pre-activation
  units one node-perturbation probe perturbs for a single example, summed over layers,
  time steps and spatial positions; it is the effective search dimension of the
  \NP{} family, as $P$ is for \SPSA{} and \RS{}.}
  \label{tab:tasks}
  \centering
  \footnotesize
  \setlength{\tabcolsep}{4pt}
  \begin{tabular}{llllrr}
    \toprule
    Key & Data & Split & Model & Params & Units \\
    \midrule
    \texttt{mlp\_mnist}       & MNIST      & 55k / 5k / 10k & MLP 784--48--32--10                          & 39,578 & 90 \\
    \texttt{conv\_cifar10}    & CIFAR-10   & 45k / 5k / 10k & conv16--conv32--conv32--fc48--10             & 39,450 & 26,682 \\
    \texttt{lstm\_ptb}        & PTB (char) & standard       & emb32--LSTM80--$V$                           & 41,810 & 25,728 \\
    \texttt{transformer\_ptb} & PTB (char) & standard       & emb48 $+$ pos64--2 blocks--$V$                  & 39,250 & 45,232 \\
    \bottomrule
  \end{tabular}
\end{table}
The character-level Penn Treebank corpus is built from the standard Mikolov
word-level files by joining lines with a newline and keeping spaces; its vocabulary has
$V=50$ symbols (5.10M / 0.40M / 0.45M characters in train / valid / test).

\paragraph{Optimization and sweep.}
Every method uses Adam with $(\beta_1,\beta_2)=(0.9,0.99)$ and learning rate
$\varepsilon$, with $\varepsilon$ also the probe scale (Section~\ref{sec:setup-rule}).
We sweep $\varepsilon\in\{0.1,\,0.01,\,0.001\}$, held fixed for the whole run. The
loss is cross-entropy: per example for classification, per sequence (mean over the 64
next-character predictions) for language modelling. Zero-order estimates use
per-example loss differences where the method allows it (the node-perturbation
family) and batch-mean losses where it does not (\SPSA{}, \RS{}).

\paragraph{Budget and evaluation.}
Each run trains for \textbf{one hour of wall-clock} on a single NVIDIA RTX~5090;
evaluation time is excluded from the clock. Validation and test loss (and accuracy or
bits-per-character) are evaluated at initialisation, every 5 minutes of training time
and at the end. A run is stopped and marked as diverged when a loss becomes NaN, or
when at two consecutive evaluations its training loss exceeds 2$\times$ its initial
value or its validation loss exceeds 2$\times$ its value at initialisation.

\paragraph{Model selection and reporting.}
For each (task, optimizer) pair we select the $\varepsilon$ with the lowest final
\emph{validation} loss and report that run's final \emph{test} loss, together with test
accuracy for classification and bits-per-character ($\text{loss}/\ln 2$) for language
modelling. The grid is 4 tasks $\times$ 8 optimizers $\times$ 3 values of
$\varepsilon$ = \textbf{96 runs}, executed 8 at a time on 8 GPUs (one run per GPU).
Every run is logged to Weights \& Biases (project \texttt{zoo-survey}, run name
\texttt{\{task\}-\{opt\}-lr\{eps\}}) and writes a JSON file with its configuration,
history and final metrics.

\paragraph{Hardware and software.}
All runs use one RTX~5090 (32\,GB) per run on a single 8-GPU node (Intel Xeon
Platinum 8490H, NVIDIA driver 580.173, CUDA 12.8, PyTorch 2.9.1, Python 3.11) rented
on vast.ai, and the code at \url{https://github.com/Fchaubard/zero-order-ablations}
(commit \texttt{db100bc} for the first ten runs and \texttt{f5424e9} for the other
86 and the supplementary runs; the latter adds the validation-loss clause of the
divergence rule, Section~\ref{sec:discussion}).

% ====================================================================================
\section{Results}
\label{sec:results}

\subsection{Main comparison}
\label{sec:main}

Table~\ref{tab:main} reports, for every task and optimizer, the final test loss of the
run selected by validation loss, with test accuracy (classification) or
bits-per-character (language modelling) as the second number. The best zero-order
method in each row is set in bold; backpropagation and forward-forward are references.
Among the probe methods, \ANP{} has the lowest final test loss on both image tasks
(0.381 on MNIST, 96.4\,\% accuracy; 1.177 on CIFAR-10, 59.8\,\%), followed by \DANP{}
(0.751 / 94.2\,\% and 1.285 / 60.4\,\%) and \INP{} (0.916 / 97.0\,\% and 1.339 /
52.7\,\%), while \INP{} is the best probe method on both sequence tasks (1.955, i.e.\
2.82 bits/char, on the LSTM and 2.151, 3.10 bits/char, on the Transformer), ahead of
\NP{} and \SPSA{} on the LSTM and of \SPSA{} and \RS{} on the Transformer. Against
backpropagation the picture splits by task family. On the sequence tasks \BP{} reaches
1.78 and 1.82 bits/char, so the best zero-order rule trails it by 1.0 and 1.3
bits/char, a factor 1.6--1.7 in test loss. On the image tasks the \emph{final} \BP{}
numbers (1.857 / 20.9\,\% on MNIST, 2.402 / 50.2\,\% on CIFAR-10) are worse than
those of the node-perturbation methods, because backpropagation completes 3--5
million steps in the hour, overfits the 40k-parameter models within minutes and then
drifts: its $\varepsilon=0.001$ MNIST run ends at 97.4\,\% accuracy but a test loss of
2.355, so validation loss selects the $\varepsilon=0.1$ run that never left
20.9\,\% (Appendix~\ref{app:grid}), and its $\varepsilon=0.001$ CIFAR-10 run was
stopped by the divergence rule at 67.4\,\% accuracy with a validation loss ten times
its initial value (Section~\ref{sec:discussion}).

Table~\ref{tab:earlystop} therefore complements the protocol's final numbers with an
early-stopping oracle: for every (task, optimizer) pair it reports the evaluation,
over all three $\varepsilon$ and all 5-minute evaluations, with the lowest validation
loss. Under this view \BP{} recovers its expected standing on the image tasks
(97.4\,\% on MNIST at minute 5; with 30-second evaluations it reaches 97.5\,\% with
test loss 0.26 after 30 seconds, and 69.7\,\% on CIFAR-10 after 30 seconds,
Section~\ref{sec:discussion}), \ANP{} wins MNIST (0.137, 96.7\,\%), \DANP{} wins
CIFAR-10 (1.085, 62.9\,\%, reached after five minutes), and the sequence-task entries
barely change because every probe method there is still improving at minute 60. The
ranking of the probe methods is stable across the two views: \INP{} has the best mean
rank (2.0 under the final-loss protocol, 1.75 under early stopping), \NP{}, \DANP{},
\SPSA{} and \ANP{} share the middle (3.25--4.25), and \RS{} and \FF{} are last.

\begin{table}[t]
  \caption{Main result: final test loss after one hour of training, with test
  accuracy (\%) for MNIST / CIFAR-10 or bits-per-character for PTB in the second row
  of each block. Each entry is the run with the lowest final validation loss over
  $\varepsilon\in\{0.1,0.01,0.001\}$. The best zero-order probe method per row is in
  bold; \BP{} is the reference and \FF{} is a local (non-probe) rule. A dagger marks
  a diverged run. For \FF{} the accuracy of label-accumulation inference is given in
  parentheses. On the image tasks the final \BP{} entries are dominated by overfitting
  (Section~\ref{sec:main}); Table~\ref{tab:earlystop} gives the early-stopped values.}
  \label{tab:main}
  \centering
  \footnotesize
  \setlength{\tabcolsep}{2.5pt}
  % PLACEHOLDER written by make_dummy_data.py -- entries are loss (acc or bpc).
% Best zero-order probe method per row should be wrapped in \textbf{}.
\begin{tabular}{l c cccccc c}
\toprule
 & BP & \multicolumn{6}{c}{Zero-order probe methods} & FF \\
\cmidrule(lr){3-8}
Task & & SPSA & RS & NP & ANP & INP & DANP & \\
\midrule
MNIST / MLP & -- (--) & -- (--) & -- (--) & -- (--) & -- (--) & -- (--) & -- (--) & -- (--) \\
CIFAR-10 / ConvNet & -- (--) & -- (--) & -- (--) & -- (--) & -- (--) & -- (--) & -- (--) & -- (--) \\
PTB / char-LSTM & -- (--) & -- (--) & -- (--) & -- (--) & -- (--) & -- (--) & -- (--) & -- (--) \\
PTB / char-Transformer & -- (--) & -- (--) & -- (--) & -- (--) & -- (--) & -- (--) & -- (--) & -- (--) \\
\bottomrule
\end{tabular}

\end{table}

\begin{table}[t]
  \caption{Early-stopping oracle: test loss and test accuracy (\%; in parentheses for
  \FF{}: label-accumulation inference) or bits per character at the evaluation with the
  lowest \emph{validation} loss over all $\varepsilon$ and all evaluation times; the
  third row of each block gives that $\varepsilon$ and the minute of training at which it
  occurred ($\varepsilon$/min; minute 0 = initialisation, i.e.\ no improvement). Best
  probe method per task in bold.}
  \label{tab:earlystop}
  \centering
  \scriptsize
  \setlength{\tabcolsep}{2pt}
  % generated by scripts/collect_results.py -- test loss / accuracy at the evaluation with the lowest validation loss over all eps and all 5-minute evaluations (early-stopping oracle); best probe method per task in bold
\begin{tabular}{ll c cccccc c}
\toprule
 & & BP & \multicolumn{6}{c}{Zero-order probe methods} & FF \\
\cmidrule(lr){4-9}
Task & Metric & & SPSA & RS & NP & ANP & INP & DANP & \\
\midrule
MNIST / MLP & test loss & 1.260 & 0.391 & 0.487 & 0.167 & \textbf{0.137} & 0.162 & 0.321 & 1.051 \\
 & acc.\ (\%) & 97.4 & 88.8 & 86.3 & 96.6 & \textbf{96.7} & 96.3 & 93.2 & 61.9 (93.2) \\
 & $\varepsilon$/min & 0.001/5 & 0.001/5 & 0.001/5 & 0.001/5 & 0.001/5 & 0.001/5 & 0.01/5 & 0.001/5 \\
\addlinespace[2pt]
CIFAR-10 / ConvNet & test loss & 1.319 & 1.642 & 1.742 & 1.434 & 1.139 & 1.341 & \textbf{1.085} & 1.819 \\
 & acc.\ (\%) & 56.0 & 40.3 & 34.9 & 48.8 & 61.7 & 53.1 & \textbf{62.9} & 34.8 (24.4) \\
 & $\varepsilon$/min & 0.01/45 & 0.001/5 & 0.001/10 & 0.001/10 & 0.001/50 & 0.001/45 & 0.001/5 & 0.001/20 \\
\addlinespace[2pt]
PTB / char-LSTM & test loss & 1.231 & 2.048 & 2.124 & 2.015 & 2.969 & \textbf{1.955} & 2.608 & 2.658 \\
 & bits/char & 1.78 & 2.96 & 3.06 & 2.91 & 4.28 & \textbf{2.82} & 3.76 & 3.84 \\
 & $\varepsilon$/min & 0.001/60 & 0.001/60 & 0.001/60 & 0.001/60 & 0.001/25 & 0.001/60 & 0.001/60 & 0.01/20 \\
\addlinespace[2pt]
PTB / char-Transformer & test loss & 1.260 & 2.181 & 2.279 & 2.342 & 3.065 & \textbf{2.151} & 2.513 & 2.925 \\
 & bits/char & 1.82 & 3.15 & 3.29 & 3.38 & 4.42 & \textbf{3.10} & 3.63 & 4.22 \\
 & $\varepsilon$/min & 0.001/55 & 0.001/60 & 0.001/60 & 0.001/35 & 0.001/10 & 0.001/60 & 0.01/50 & 0.01/25 \\
\bottomrule
\end{tabular}

\end{table}

\subsection{Training curves}
\label{sec:curves}

Figure~\ref{fig:curves} shows test loss against training wall-clock for the selected
$\varepsilon$ of every method on every task. Plotting against time rather than steps
is the fair comparison under our protocol: methods that spend 16 passes per step take
proportionally fewer steps per hour, and sequence models pay for every pass with $T=64$
sequential time steps.
The curves have three shapes. On MNIST every probe method reaches its minimum at the
first evaluation (minute 5: \ANP{} 0.137, \INP{} 0.162, \NP{} 0.167) and the test loss
then rises for the rest of the hour while the accuracy barely moves (\NP{}: 96.6\,\%
to 96.8\,\% as the loss climbs from 0.167 to 2.141), i.e.\ the networks become
over-confident rather than wrong; the same happens to \BP{} within its first minute
(Section~\ref{sec:discussion}). On CIFAR-10 the methods separate: \DANP{} drops to
1.09 within five minutes and then creeps up, \ANP{} and \INP{} improve slowly for the
whole hour (\ANP{} reaches its minimum at minute 50), \NP{} bottoms out at minute 10
and drifts, and \SPSA{}, \RS{} and \FF{} stay near the chance loss of 2.30. On both
sequence tasks every probe method is still improving at minute 60, slowly, on a
plateau roughly one bit per character above \BP{}, which is itself still improving;
the Transformer panel shows the spikes of \ANP{}, whose updates are erratic on
sequence models (Section~\ref{sec:analysis}). Appendix~\ref{app:curves} shows the
curves of all 96 runs: at $\varepsilon=0.1$ most probe methods diverge within the
first two evaluations, and at $\varepsilon=0.01$ several reach a good minimum early
and then blow up (\INP{} on MNIST, \NP{} on CIFAR-10, \ANP{} and \DANP{} on the LSTM).

\begin{figure}[t]
  \centering
  % Figure 1: 2x2 grid of test loss vs training time, one panel per task.
% Data: data/curve_{task}.csv with columns time,bp,spsa,rs,np,anp,inp,danp,ff
% The legend is written to the name "zoolegend" (first panel) and placed with
% \ref{zoolegend} in main.tex; the same legend is reused by the other figures.
\begin{tikzpicture}
  \begin{groupplot}[
      group style={group size=2 by 2, horizontal sep=1.5cm, vertical sep=1.6cm},
      width=0.5\linewidth, height=4.4cm,
      cycle list name=zoocycle,
      ymode=log,
      xmin=0, xmax=60,
      xlabel={training time (min)}, ylabel={test loss},
      tick label style={font=\scriptsize},
      label style={font=\scriptsize},
      title style={font=\small, yshift=-1ex},
      every axis plot/.append style={line width=0.7pt},
      grid=major, grid style={gray!25},
    ]
    \nextgroupplot[title={MNIST / MLP},
      legend to name=zoolegend, legend columns=8,
      legend style={font=\scriptsize, draw=none, /tikz/every even column/.append style={column sep=0.5em}},
      legend entries={BP, SPSA, RS, NP, ANP, INP, DANP, FF}]
    \zooplots{time}{data/curve_mlp_mnist.csv}

    \nextgroupplot[title={CIFAR-10 / ConvNet}]
    \zooplots{time}{data/curve_conv_cifar10.csv}

    \nextgroupplot[title={PTB / char-LSTM}]
    \zooplots{time}{data/curve_lstm_ptb.csv}

    \nextgroupplot[title={PTB / char-Transformer}]
    \zooplots{time}{data/curve_transformer_ptb.csv}
  \end{groupplot}
\end{tikzpicture}

  \\[0.5ex]
  \ref{zoolegend}
  \caption{Test loss versus training time (minutes, evaluation time excluded) for the
  $\varepsilon$ selected by validation loss, one line per optimizer. Log-scale
  $y$-axis, capped at 10. \BP{} is the black dashed reference; on MNIST its selected
  run is $\varepsilon=0.1$, which never leaves 20.9\,\% accuracy, because the
  $\varepsilon=0.001$ run overfits to a higher final loss at 97.4\,\% accuracy
  (Appendix~\ref{app:curves}).}
  \label{fig:curves}
\end{figure}

\subsection{Sensitivity to the probe scale}
\label{sec:eps}

Because $\varepsilon$ is both the learning rate and the probe scale, its sweep probes
two things at once: how large a step each method can tolerate and how small a probe
each method can resolve. Figure~\ref{fig:eps} plots the final test loss of every run
against $\varepsilon$.
The sweep is decisive. At $\varepsilon=0.1$, 15 of the 24 probe-method runs diverge
(all four \DANP{} runs, three each of \NP{}, \ANP{} and \RS{}, two of \SPSA{}, none
of \INP{}) and the nine survivors sit at chance or above the uniform loss; \BP{}
still trains the language models at this step size (2.57 and 2.65 bits/char) but
collapses on the image tasks, and \FF{} never diverges but its readout is at chance.
At $\varepsilon=0.01$ seven probe runs diverge, five of them node-perturbation runs,
and \BP{} itself is stopped on MNIST. At $\varepsilon=0.001$ no probe run diverges,
25 of the 32 (task, method) pairs select this value, and only \SPSA{} and \RS{} on
CIFAR-10 (12.0 and 9.4\,\%) and the \FF{} readout on MNIST fail to move. Two patterns
stand out. First, the node-perturbation family tolerates less noise than the
weight-space methods: it accounts for 10 of the 15 divergences at $\varepsilon=0.1$
and 5 of 7 at $\varepsilon=0.01$, consistent with the node noise entering every
layer's pre-activations directly, and \DANP{} is the most fragile of all
(Section~\ref{sec:discussion}). Second, \INP{} is the only probe method with no
divergence at $\varepsilon=0.1$; its single flag (MNIST, $\varepsilon=0.01$) was
raised at the final evaluation, after an hour of oscillation around 70\,\% accuracy.

\begin{figure}[t]
  \centering
  % Figure 2: epsilon sensitivity -- final test loss vs epsilon, one panel per task,
% one line per optimizer.  Data: data/lr_{task}.csv with columns
% eps,bp,spsa,rs,np,anp,inp,danp,ff  (three rows: eps = 0.1, 0.01, 0.001).
\begin{tikzpicture}
  \begin{groupplot}[
      group style={group size=2 by 2, horizontal sep=1.5cm, vertical sep=1.6cm},
      width=0.5\linewidth, height=4.2cm,
      cycle list name=zoocycle,
      xmode=log, ymode=log, log ticks with fixed point, ymax=10,
      log basis x=10,
      xtick={0.001, 0.01, 0.1},
      xticklabels={$10^{-3}$, $10^{-2}$, $10^{-1}$},
      xmin=0.0005, xmax=0.2,
      xlabel={$\varepsilon$ (learning rate $=$ probe scale)}, ylabel={final test loss},
      tick label style={font=\scriptsize},
      label style={font=\scriptsize},
      title style={font=\small, yshift=-1ex},
      every axis plot/.append style={line width=0.7pt}, unbounded coords=jump,
      grid=major, grid style={gray!25},
    ]
    \nextgroupplot[title={MNIST / MLP}]
    \zooplots{eps}{data/lr_mlp_mnist.csv}

    \nextgroupplot[title={CIFAR-10 / ConvNet}]
    \zooplots{eps}{data/lr_conv_cifar10.csv}

    \nextgroupplot[title={PTB / char-LSTM}]
    \zooplots{eps}{data/lr_lstm_ptb.csv}

    \nextgroupplot[title={PTB / char-Transformer}]
    \zooplots{eps}{data/lr_transformer_ptb.csv}
  \end{groupplot}
\end{tikzpicture}

  \\[0.5ex]
  \ref{zoolegend}
  \caption{Final test loss versus $\varepsilon$ (learning rate $=$ probe scale) for
  every optimizer and task; all 96 runs. Diverged runs are plotted at their last
  finite evaluation.}
  \label{fig:eps}
\end{figure}

\subsection{Throughput}
\label{sec:throughput}

Table~\ref{tab:throughput} lists the number of optimizer steps and forward passes
completed per hour by every method on every task. Since the per-step budget is 16
passes for the probe methods, 2 for \FF{} and 1 (plus a backward) for \BP{}, the
step rate is the quantity that differs across methods, while the pass rate shows how
efficiently each implementation uses the GPU (batched probes, in-place perturbation,
sequential LSTM time steps).
Sixteen batched probe passes cost 4.3--4.7 backpropagation steps on every task
(\SPSA{} completes 1.09M, 0.62M, 34k and 219k steps per hour on the four tasks
against 4.94M, 2.93M, 148k and 1.02M for \BP{}), so one probe pass costs about 0.28
of a forward-plus-backward and the probe methods see 3.4--3.7 times more forward
passes per hour than \BP{}: the much lower step count is the price of the estimator,
not of the implementation. Within the probe family, \NP{} is 1.5--1.7 times slower
than \SPSA{} because its clean pass builds a local graph per layer and runs one short
backward per layer; \ANP{} adds 20--36\,\% for the per-example activity norms;
\DANP{} adds another 16--26\,\% for the decorrelation update and is 1.4--1.7 times
slower than \NP{} (2.9 times slower than \SPSA{} on the MLP); \INP{} is as fast as
\NP{} on the MLP and up to 34\,\% faster on the Transformer because each noisy pass
perturbs one layer only. \FF{} takes 2.7 times fewer steps than \BP{} on the image
models despite its two-pass budget, because it forwards a triple batch (positive,
negative and neutral) and runs one backward per block, but matches \BP{} on the
LSTM, whose cost is dominated by the sequential time steps. The LSTM is the slowest
model for every method: at $T=64$ sequential steps per pass, \SPSA{} completes 32
times fewer steps per hour than on the MLP, and the probe methods manage only
17k--35k steps in the hour.

\begin{table}[t]
  \caption{Throughput: optimizer steps per hour and forward passes per hour (in
  thousands) for every method and task, measured on one RTX~5090 at the selected
  $\varepsilon$.}
  \label{tab:throughput}
  \centering
  \footnotesize
  \setlength{\tabcolsep}{2pt}
  % generated by scripts/collect_stage2.py -- optimizer steps per hour and forward passes per hour (thousands), mean over seeds of the selected runs
\begin{tabular}{l cc cc cc cc}
\toprule
 & \multicolumn{2}{c}{MNIST} & \multicolumn{2}{c}{CIFAR-10} & \multicolumn{2}{c}{LSTM} & \multicolumn{2}{c}{Transformer} \\
\cmidrule(lr){2-3}\cmidrule(lr){4-5}\cmidrule(lr){6-7}\cmidrule(lr){8-9}
Method & steps/h & passes/h & steps/h & passes/h & steps/h & passes/h & steps/h & passes/h \\
\midrule
\BP{} & 4278.4 & 4278.4 & 2023.8 & 2023.8 & 147.5 & 147.5 & 1021.0 & 1021.0 \\
\SPSA{} & 1002.7 & 16043.3 & 590.4 & 9446.4 & 34.3 & 548.1 & 219.1 & 3506.3 \\
\SPSAb{} & 850.3 & 14455.1 & 511.2 & 8690.2 & 33.6 & 571.8 & 227.2 & 3862.4 \\
\RS{} & 1075.8 & 17212.5 & 609.0 & 9743.7 & 34.5 & 552.7 & 229.1 & 3665.6 \\
\NP{} & 587.4 & 9397.7 & 235.9 & 3773.9 & 23.6 & 377.4 & 74.7 & 1195.1 \\
\ANP{} & 493.3 & 7892.1 & 177.6 & 2842.0 & 19.8 & 317.2 & 62.5 & 999.6 \\
\ANPb{} & 502.9 & 8549.5 & 351.9 & 5983.1 & 23.5 & 399.0 & 73.6 & 1251.2 \\
\INP{} & 598.7 & 9579.0 & 283.6 & 4537.6 & 24.9 & 398.9 & 168.0 & 2688.0 \\
\DANP{} & 430.3 & 6885.3 & 243.3 & 3892.3 & 17.0 & 272.5 & 86.8 & 1388.1 \\
\FF{} & 1685.1 & 5055.3 & 941.0 & 2823.0 & 139.7 & 279.4 & 271.0 & 542.0 \\
\bottomrule
\end{tabular}

\end{table}

\subsection{Analysis}
\label{sec:analysis}

\paragraph{Which zero-order method wins?}
\INP{} is the answer if a single rule must be chosen: it is the best probe method on
both sequence tasks, second or third in test loss on both image tasks with the
highest final MNIST accuracy of any probe method (97.0\,\%), has the best mean rank
under both views of the results (2.0 and 1.75 of 7) and is the only probe method that
survives $\varepsilon=0.1$ on every task. \ANP{} is the best rule on the image tasks
under the protocol (0.381 and 1.177, against \INP{}'s 0.916 and 1.339) and \DANP{} is
the fastest to a good CIFAR-10 model (62.9\,\% after five minutes), but both fail on
the sequence models, where \ANP{} is the worst probe method (4.66 and 4.92 bits/char)
and \DANP{} fifth. The gap to backpropagation depends on the task family. On the
image tasks it is small at the best evaluation (96.7 versus 97.5\,\% on MNIST, 62.9
versus 69.7\,\% on CIFAR-10) and the protocol's final numbers even favour the probe
methods, which take 7--13 times fewer steps than \BP{} and therefore do not overfit
within the hour. On the sequence tasks the best probe method ends 1.0 (LSTM) and 1.3
(Transformer) bits per character behind \BP{}, a factor 1.6--1.7 in test loss, and is
still on a slowly descending plateau at minute 60.

\paragraph{Weight probes versus node probes.}
\SPSA{} and \RS{} perturb all $P\approx 40$k weights at once and obtain one scalar
loss difference per probe, so the variance of their estimate grows with $P$ and 16
passes yield a direction whose cosine with the gradient is of order
$\sqrt{16/P}\approx 0.02$. The node-perturbation family obtains a loss difference
\emph{per example} (128 or 64 per pass) and assigns credit per unit, so its effective
dimension is the number of perturbed units (Table~\ref{tab:tasks}): 90 on the MLP
against 39,578 weights. The results track this ratio. On MNIST the node methods reach
96--97\,\% while \SPSA{} and \RS{} reach 72 and 62\,\% (89 and 86\,\% at their best
evaluation). On the ConvNet the 26,682 perturbed unit-positions are of the same order
as the 39,450 weights and plain \NP{} drops to 41.9\,\%; we measured the cosine
between the \NP{} estimate and the true gradient on this network and found it at 0.0
for the first two convolutional layers with 480 probes (0.55 for the third, 0.94 and
0.99 for the dense layers), rising to 0.59 only with 7,680 probes, so with 15 probes
per step the early layers receive essentially no signal and Adam, which normalises
every coordinate, turns the noise into a random walk of fixed size $\varepsilon$. On
the sequence models the number of unit-positions (25,728 and 45,232) exceeds $P$, and
the two families end within 0.1--0.3 bits/char of each other (\SPSA{} 2.96 and 3.15,
\INP{} 2.82 and 3.10), with \RS{} close behind. The weight-space methods are,
however, markedly more tolerant of large $\varepsilon$ and 1.5--2.9 times faster per
step.

\paragraph{The effect of activity-based, iterative and decorrelated variants.}
\ANP{}'s activity-based attribution clearly helps on the feed-forward image models:
it beats \NP{} by a wide margin in final loss (0.381 versus 2.141 on MNIST, 1.177
versus 1.613 on CIFAR-10) and in accuracy on CIFAR-10 (59.8 versus 41.9\,\%), in line
with \citet{dalm2023}. On the recurrent and attention models it is a liability:
\ANP{} diverges at $\varepsilon\ge 0.01$ on both and barely leaves the unigram level
at $\varepsilon=0.001$ (4.66 and 4.92 bits/char against 2.91 and 3.40 for \NP{}). The
estimator itself is sound there---on the LSTM its cosine with the true gradient is
0.45--0.82 with 480 probes, close to \NP{}'s 0.55--0.95---but its normalisation by
the squared activity difference, summed over 25k--45k unit-positions whose
perturbations propagate through shared weights, produces updates of erratic magnitude
(the spikes in Figure~\ref{fig:curves}). \INP{}'s one-layer-at-a-time attribution has
only $15/L$ draws per layer (5, 3, 5 and 1 on the four models), yet removes all
cross-layer noise from each estimate, and this trade is worth it on every task:
\INP{} is never worse than third and is the best rule on both sequence models,
including the Transformer with a single draw per layer. Decorrelation helps where
\citet{dalm2023} reported it and not elsewhere: \DANP{} is the fastest method on
CIFAR-10 (1.085 at minute 5, a level \ANP{} needs 50 minutes to reach) and improves
on \ANP{} on both sequence models (3.76 and 3.64 versus 4.66 and 4.92 bits/char), but
ends behind \ANP{} on CIFAR-10 after the full hour (1.285 versus 1.177) and clearly
behind on MNIST (0.751 versus 0.381, 94.2 versus 96.4\,\%); it is also the most
fragile rule (five of its twelve runs diverged, including NaN within 21 steps at
$\varepsilon=0.1$ on both image models, Section~\ref{sec:discussion}) and the slowest
(Table~\ref{tab:throughput}). Decorrelation is not applied to the embedding layers,
so on the sequence models it acts on the recurrent / attention and output layers only.

\paragraph{Sequence models.}
The sequence models are the hardest setting for every zero-order rule and the most
expensive: a pass through the LSTM is 64 sequential time steps, so the probe methods
complete only 17k--35k steps in the hour (Table~\ref{tab:throughput}). On the test
split, a uniform model scores 5.64 bits/char, a unigram 4.32, a bigram 3.31 and an
add-$0.1$ trigram 2.61. Every probe method beats the unigram on both models; on the
LSTM, \INP{}, \NP{}, \SPSA{} and \RS{} (2.82--3.06) beat the bigram and sit between
bigram and trigram, \DANP{} (3.76) and \FF{} (3.88) do not, and \ANP{} (4.66) is near
the unigram; on the Transformer only \INP{}, \SPSA{} and \RS{} (3.10, 3.15 and 3.29)
edge past the bigram. \BP{} reaches 1.78 and 1.82 bits/char, well beyond the trigram,
so after one hour the zero-order rules have learned roughly bigram statistics while
backpropagation has learned a real character model. All probe curves are still
descending at minute 60 (Figure~\ref{fig:curves}), so part of this gap is budget
rather than estimator, but at 17k--35k steps per hour closing it would take days.
\FF{} with corrupted-sequence negatives trains fast (140k and 778k steps per hour)
and its goodness does separate real from corrupted text, but the normalised features
support only a weak next-character readout (3.88 and 4.27 bits/char, between unigram
and bigram).

\paragraph{Forward-forward.}
\FF{} is not a probe method: it takes 1.8M steps per hour on the MLP (2.9 times
\NP{}'s count, 37\,\% of \BP{}'s), uses exact local gradients within each block and
never diverged at any $\varepsilon$. Its test loss, however, is set by the linear
readout on the neutral-label features of the last block, and that readout is poor on
the image tasks: 19.1\,\% on MNIST and 33.0\,\% on CIFAR-10, while Hinton's
label-accumulation inference on the same trained blocks gives 96.0\,\% and 32.7\,\%.
On MNIST the readout even degrades during training (61.9\,\% at minute 5, 19.1\,\%
at minute 60) and its own training cross-entropy rises towards $\ln 10$, i.e.\ the
neutral-label activity of a 32-unit block loses class information as the goodness
objective is optimised; the readout follows Hinton's recipe, but that recipe assumed
layers of 2,000 units. Label-accumulation accuracy is also the most
$\varepsilon$-robust number in the study (92.2, 95.0 and 96.0\,\% on MNIST at
$\varepsilon=0.1$, $0.01$ and $0.001$). On the sequence models \FF{} is the second
worst rule (3.88 and 4.27 bits/char), ahead only of \ANP{}.

% ====================================================================================
\section{Discussion and Limitations}
\label{sec:discussion}

\paragraph{Final loss, overfitting and the divergence rule.}
The protocol reports the loss after a fixed hour and selects $\varepsilon$ by final
validation loss. For backpropagation, which completes 3--5 million steps on the
image models, this measures overfitting rather than optimisation: with 30-second
evaluations (four supplementary 10-minute runs, released with the logs) \BP{}
reaches 97.5\,\% on MNIST and 69.7\,\% on CIFAR-10 after 30 seconds, and its test
loss then grows while the accuracy holds. The divergence rule, which stops a run
whose validation loss exceeds twice its initial value at two consecutive
evaluations, was added to catch probe runs whose training loss is masked by the probe
noise, and it also stopped three overfitting runs (\BP{} on CIFAR-10 at
$\varepsilon=0.001$ and on MNIST at $\varepsilon=0.01$, \INP{} on MNIST at
$\varepsilon=0.01$); the first ten runs of the sweep, all at $\varepsilon=0.1$,
started before this clause existed, and none of them would have been stopped earlier
by it. We keep the protocol numbers in Table~\ref{tab:main} because they are what the
pre-registered procedure produces, and report the early-stopping oracle
(Table~\ref{tab:earlystop}) alongside; weight decay, early stopping or a smaller
learning rate for \BP{} would make the reference stronger on the image tasks without
changing the ranking of the zero-order rules.

\paragraph{\DANP{}'s decorrelation rate.}
The decorrelation matrices are updated with a fixed rate $\eta_R=10^{-3}$ that is
not normalised by the input scale. At $\varepsilon=0.1$ a single Adam step inflates
the ReLU activities of the hidden layers by an order of magnitude, the update
$\bR_l \leftarrow \bR_l - \eta_R(\bC_l - \mathrm{diag}\,\bC_l)\bR_l$ then grows
$\bR_l$ tenfold per step, and the run is NaN within 21 steps on both image models
(the input-layer $\bR$, whose inputs are bounded pixels, stays stable). A rate tied
to the activity scale would remove this failure mode but is not part of the
published rule.

\paragraph{Single seed.}
Each of the 96 configurations is run once. The reported differences are therefore
single draws of a random process (initialisation, data order, probes); we do not
report error bars, and small differences in Table~\ref{tab:main} should not be read as
rankings. The one-hour-per-run protocol was chosen so
that the full grid runs in about 11 hours of wall-clock on 8 GPUs (100 GPU-hours); repeating it with several seeds is the most
important follow-up.

\paragraph{One-hour budget and 40k parameters.}
The budget is deliberately small and identical for all methods, which favours methods
that make progress quickly over methods that converge slowly but further. The
$\approx$40k-parameter models are small by modern standards; the variance of
weight-space estimators grows with the number of parameters, so the relative standing
of \SPSA{} and \RS{} against the node-perturbation family may change at larger scale,
and none of our conclusions should be extrapolated to models with orders of magnitude
more parameters without new experiments.

\paragraph{$\varepsilon$ tied to the learning rate.}
Using one scalar for the probe scale and the Adam learning rate removes a dimension
from the sweep but also removes a degree of freedom: a method that would benefit from a
small probe and a large step (or vice versa) is handicapped. The three-point sweep is
coarse, and the best $\varepsilon$ of a method may lie between grid points.

\paragraph{Wall-clock depends on implementation efficiency.}
Comparing against training time rewards implementations that batch probes well and
penalises methods whose natural implementation is sequential (iterative NP, the
decorrelation update, LSTM time steps). We report forward passes per hour
(Table~\ref{tab:throughput}) so that readers can re-weight the comparison by passes
instead of time, but the ranking under a pass budget could differ.

\paragraph{\FF{} and \BP{} are not probe methods.}
Forward-forward uses local gradients within each layer and backpropagation uses the
exact gradient; neither obeys the 16-pass budget, and both are included only as
reference points for a local rule and for the ceiling.

\paragraph{Character-level PTB.}
Language modelling is done at the character level because a word-level embedding of
the 10k-word PTB vocabulary would by itself exceed the parameter budget. Character-level
results are not comparable to word-level perplexities in the literature.

% ====================================================================================
\section{Conclusion}
\label{sec:conclusion}

We compared eight training rules---backpropagation, SPSA, random search, four
node-perturbation variants and forward-forward---on four $\approx$40k-parameter
architectures under one fully specified protocol: the same 16 forward passes per step,
the same hour of training, the same Adam update and a single swept scalar that is both
learning rate and probe scale.
If one zero-order rule must be chosen, it is iterative node perturbation: it is the
best probe method on both sequence models, within a point of the best on both image
models, the most robust to the probe scale and the only variant with no divergence at
$\varepsilon=0.1$; activity-based and decorrelated node perturbation win the image
tasks but fail on sequences, and weight-space perturbation is competitive only where
the number of perturbed units approaches the number of weights. At their best
evaluation the node-perturbation rules are within a point of backpropagation on MNIST
and seven points on CIFAR-10, but on the character-level LSTM and Transformer the gap
is still about one bit per character after an hour, with every zero-order curve on a
slowly descending plateau.
The protocol, code and all 96 run logs are public so that new zero-order rules can be
added to the grid at the cost of 12 GPU-hours.

% ====================================================================================
\bibliographystyle{plainnat}
\bibliography{refs}

% ====================================================================================
\appendix

\section{Full 96-run grid}
\label{app:grid}

Table~\ref{tab:grid} lists the final validation and test loss of every run, for all
three values of $\varepsilon$; the run selected for Table~\ref{tab:main} in each row is
the one with the lowest validation loss.

\begin{table}[ht]
  \caption{Final validation / test loss of all 96 runs. For each (task, optimizer)
  row the $\varepsilon$ with the lowest validation loss (in bold) is the run reported
  in Table~\ref{tab:main}. A dagger marks a diverged run (last finite evaluation
  shown).}
  \label{tab:grid}
  \centering
  \footnotesize
  \setlength{\tabcolsep}{4pt}
  % generated by scripts/collect_results.py -- validation / test loss for every eps; selected eps in bold; $\dagger$ = diverged (stopped early)
\begin{tabular}{ll ccc}
\toprule
 & & \multicolumn{3}{c}{validation / test loss} \\
\cmidrule(lr){3-5}
Task & Optimizer & $\varepsilon=0.1$ & $\varepsilon=0.01$ & $\varepsilon=0.001$ \\
\midrule
MNIST / MLP & BP & \textbf{1.826 / 1.857} & 51.107 / 177.402$^\dagger$ & 2.239 / 2.355 \\
 & SPSA & 957.457 / 2.302$^\dagger$ & 2.193 / 5.502 & \textbf{0.917 / 1.039} \\
 & RS & 52.924 / 170.573$^\dagger$ & 8.157 / 4.833 & \textbf{1.238 / 1.412} \\
 & NP & 2.301 / 2.301$^\dagger$ & 2.303 / 2.302 & \textbf{1.985 / 2.141} \\
 & ANP & 2.323 / 2.319 & 0.519 / 0.619 & \textbf{0.363 / 0.381} \\
 & INP & 2.325 / 2.321 & 8.917 / 21.010$^\dagger$ & \textbf{0.935 / 0.916} \\
 & DANP & -- / --$^\dagger$ & 1.033 / 1.190 & \textbf{0.643 / 0.751} \\
 & FF & 2.288 / 2.288 & 2.243 / 2.241 & \textbf{2.123 / 2.127} \\
\midrule
CIFAR-10 / ConvNet & BP & 2.305 / 2.305 & \textbf{1.357 / 2.402} & 23.735 / 24.475$^\dagger$ \\
 & SPSA & 7.292 / 7.225$^\dagger$ & 22.978 / 2.313$^\dagger$ & \textbf{2.300 / 2.468} \\
 & RS & \textbf{2.395 / 2.406} & 197.361 / 38.724$^\dagger$ & 2.401 / 2.492 \\
 & NP & $1.9{\times}10^{6}$ / $1.3{\times}10^{6}$$^\dagger$ & 314.593 / 326.902$^\dagger$ & \textbf{1.597 / 1.613} \\
 & ANP & 23.474 / 23.143$^\dagger$ & 2.332 / 2.336 & \textbf{1.143 / 1.177} \\
 & INP & 2.327 / 2.323 & 2.303 / 2.303 & \textbf{1.321 / 1.339} \\
 & DANP & -- / --$^\dagger$ & \textbf{1.242 / 1.285} & 1.289 / 1.371 \\
 & FF & 2.311 / 2.311 & 2.081 / 2.084 & \textbf{1.880 / 1.886} \\
\midrule
PTB / char-LSTM & BP & 1.858 / 1.835 & 1.371 / 1.340 & \textbf{1.268 / 1.231} \\
 & SPSA & 3.991 / 3.965 & 2.125 / 2.107 & \textbf{2.066 / 2.048} \\
 & RS & 15.802 / 15.689$^\dagger$ & 2.330 / 2.310 & \textbf{2.141 / 2.124} \\
 & NP & 23.036 / 22.815$^\dagger$ & 2.380 / 2.356 & \textbf{2.033 / 2.015} \\
 & ANP & 195.280 / 192.326$^\dagger$ & 9.592 / 9.603$^\dagger$ & \textbf{3.242 / 3.232} \\
 & INP & 2.885 / 2.866 & 2.007 / 1.989 & \textbf{1.971 / 1.955} \\
 & DANP & 271.941 / 271.567$^\dagger$ & 13.384 / 13.251$^\dagger$ & \textbf{2.627 / 2.608} \\
 & FF & 2.862 / 2.859 & \textbf{2.699 / 2.688} & 2.789 / 2.776 \\
\midrule
PTB / char-Transformer & BP & 1.802 / 1.780 & 1.351 / 1.316 & \textbf{1.300 / 1.262} \\
 & SPSA & 6.865 / 6.893 & 2.557 / 2.539 & \textbf{2.201 / 2.181} \\
 & RS & 9.453 / 9.538$^\dagger$ & 2.745 / 2.726 & \textbf{2.297 / 2.279} \\
 & NP & 6.617 / 6.746 & 3.206 / 3.180 & \textbf{2.374 / 2.357} \\
 & ANP & 305.304 / 306.216$^\dagger$ & 27.783 / 27.601$^\dagger$ & \textbf{3.433 / 3.412} \\
 & INP & 5.921 / 5.880 & 2.651 / 2.638 & \textbf{2.170 / 2.151} \\
 & DANP & 26.685 / 26.551$^\dagger$ & \textbf{2.543 / 2.525} & 2.564 / 2.554 \\
 & FF & 5.707 / 5.694 & \textbf{2.958 / 2.959} & 2.976 / 2.983 \\
\bottomrule
\end{tabular}

\end{table}


\section{Implementation details and hyper-parameters}
\label{app:impl}

\paragraph{Perturbable layers.}
Every trainable layer of every model is a perturbable linear, convolutional or
embedding layer. A global context selects what a layer does on its forward pass:
(i)~a plain differentiable forward (used by backpropagation, by the forward-forward
local losses and by evaluation); (ii)~a \emph{clean} forward that detaches the layer
input and records the input $\bh_{l-1}$ and pre-activation $\ba_l$, keeping a local
autograd graph to the layer's own parameters only; or (iii)~a graph-free \emph{noisy}
forward that adds $\sigma\bxi_l$ to the pre-activation of every layer (or of one
selected layer for \INP{}) and records $\bxi_l$ and $\tilde\ba_l$. Given a
per-unit credit tensor $\mathbf{G}_l$ (for \NP{}, $\dL_n\bxi_{l,n}/\sigma$), the weight
update $\mathbf{G}_l^{\top}\bh_{l-1}$ and its convolutional and embedding analogues are
obtained by a one-layer backward through the recorded local graph, so that no error
is propagated through any nonlinearity. A layer that is called several times in one
forward (the LSTM gate layer, once per time step) accumulates one record per call and
the credit is summed over calls.

\paragraph{Decorrelation.}
For \DANP{}, each linear and convolutional layer owns a square matrix $\bR_l$
(initialised to the identity, stored as a buffer, not an optimizer parameter) that
multiplies the layer input; for convolutions the decorrelated vector is the unfolded
$C_{\text{in}}k^{2}$ patch. $\bR_l$ is updated once per step with
Eq.~\eqref{eq:decor} from the inputs recorded on the clean pass with $\eta_R=10^{-3}$, from the very first step (no warm-up). Embedding
layers are not decorrelated.

\paragraph{Language-model specifics.}
The LSTM is written with a single perturbable layer for the four gate pre-activations so
that the gates are the perturbable nodes; the forget-gate bias is initialised to 1. The
Transformer uses pre-LN blocks with parameter-free layer normalisation, a learned
positional embedding over $T=64$ positions, and a final parameter-free layer
normalisation before the readout. The per-sequence loss is the mean next-character
cross-entropy over the 64 positions.

\paragraph{Forward-forward specifics.}
One global Adam instance updates all parameters; since every block's goodness loss is
computed on the detached, normalised output of the previous block, each parameter
receives exactly one local gradient per step (the token and position embeddings are
trained by the first block's goodness). The label overlay follows the reference
implementation of \citet{hinton2022}: the first 10 entries of the flattened input are
zeroed and the entry of the label is set to the maximum input value (for CIFAR-10 these
are the first 10 pixels of the first row of channel 0); a wrong label is drawn uniformly
from the nine other classes; the neutral label sets all 10 entries to $0.1$ times the
maximum. For PTB, 30\,\% of the positions of each sequence are replaced by uniformly
random characters, resampled every step. The readout is trained every step by
cross-entropy on the normalised, detached top-layer activity (neutral-label pass for
images; the real sequences at every position for text). Goodness threshold $\theta=2$
for all blocks except the LSTM block ($\theta=0.25$, Section~\ref{sec:methods}).

\paragraph{Evaluation.}
Validation and test metrics are computed with the plain forward of the full
validation / test set at initialisation, every 5 minutes of training time (the clock
is paused during evaluation) and at 60 minutes. For \FF{}, evaluation uses the neutral-label pass and
the trained readout.

\paragraph{Hyper-parameters.}
Table~\ref{tab:hparams} summarises every hyper-parameter; none is tuned per method
beyond the $\varepsilon$ sweep.

\begin{table}[ht]
  \caption{Hyper-parameters shared by all runs.}
  \label{tab:hparams}
  \centering
  \footnotesize
  \begin{tabular}{l p{7.6cm}}
    \toprule
    Optimizer & Adam, $(\beta_1,\beta_2)=(0.9,0.99)$, $\epsilon_{\text{Adam}}=10^{-8}$, no schedule, no weight decay \\
    Learning rate $=$ probe scale $\varepsilon$ & $\{0.1, 0.01, 0.001\}$ \\
    Probe distribution & $\mathcal{N}(\mathbf{0},\bI)$, scaled by $\varepsilon$ \\
    Forward passes per step & 16 (probe methods), 2 (\FF{}), 1 $+$ 1 backward (\BP{}) \\
    \SPSA{} antithetic pairs / \RS{} candidates & 8 / 15 \\
    Batch size & 128 images; 64 sequences $\times$ 64 characters \\
    Initialisation & Xavier-uniform, zero biases (LSTM forget gate $=1$) \\
    \FF{} threshold $\theta$ / PTB corruption & 2 (0.25 for the LSTM block) / 30\,\% of characters \\
    Decorrelation rate $\eta_R$ (\DANP{}) & $10^{-3}$, fixed \\
    Training budget & 60 min wall-clock per run, evaluation excluded \\
    Evaluation interval & every 5 min of training time, and at the end \\
    Divergence rule & NaN, or training loss $>2\times$ its initial value, or validation loss $>2\times$ its step-0 value, at two consecutive evaluations \\
    Seeds & 1 per configuration \\
    \bottomrule
  \end{tabular}
\end{table}

\paragraph{Data and licences.}
MNIST \citep{lecun1998}, CIFAR-10 \citep{krizhevsky2009} and the character-level Penn
Treebank files \citep{marcus1993,mikolov2012} are used as distributed by their
authors. MNIST is distributed under the Creative Commons Attribution-Share Alike 3.0
licence; CIFAR-10 is released by its authors for research use without a formal
licence; the character-level PTB files are the widely redistributed Mikolov
preprocessing of the Penn Treebank (LDC99T42) and are used for research only. Our code
is released under the MIT licence.

\section{Per-task training curves for all \texorpdfstring{$\varepsilon$}{epsilon}}
\label{app:curves}

Figure~\ref{fig:curves-all} shows the test-loss curves of all 96 runs, one row per
task and one column per $\varepsilon$.

\begin{figure}[ht]
  \centering
  % Appendix C: test loss vs training time for every run.
% Rows = tasks, columns = epsilon in {0.1, 0.01, 0.001}.
% Data: data/curve_{task}_lr{eps}.csv with columns time,bp,spsa,rs,np,anp,inp,danp,ff
\begin{tikzpicture}
  \begin{groupplot}[
      group style={group size=3 by 4, horizontal sep=1.0cm, vertical sep=1.3cm,
                   xlabels at=edge bottom, ylabels at=edge left},
      width=0.33\linewidth, height=3.9cm,
      cycle list name=zoocycle,
      ymode=log,
      xmin=0, xmax=60,
      xlabel={training time (min)}, ylabel={test loss},
      tick label style={font=\scriptsize},
      label style={font=\scriptsize},
      title style={font=\scriptsize, yshift=-1ex},
      every axis plot/.append style={line width=0.6pt},
      grid=major, grid style={gray!25},
    ]
    % ---- MNIST / MLP
    \nextgroupplot[title={MNIST / MLP, $\varepsilon=0.1$}]
    \zooplots{time}{data/curve_mlp_mnist_lr0.1.csv}
    \nextgroupplot[title={MNIST / MLP, $\varepsilon=0.01$}]
    \zooplots{time}{data/curve_mlp_mnist_lr0.01.csv}
    \nextgroupplot[title={MNIST / MLP, $\varepsilon=0.001$}]
    \zooplots{time}{data/curve_mlp_mnist_lr0.001.csv}
    % ---- CIFAR-10 / ConvNet
    \nextgroupplot[title={CIFAR-10 / ConvNet, $\varepsilon=0.1$}]
    \zooplots{time}{data/curve_conv_cifar10_lr0.1.csv}
    \nextgroupplot[title={CIFAR-10 / ConvNet, $\varepsilon=0.01$}]
    \zooplots{time}{data/curve_conv_cifar10_lr0.01.csv}
    \nextgroupplot[title={CIFAR-10 / ConvNet, $\varepsilon=0.001$}]
    \zooplots{time}{data/curve_conv_cifar10_lr0.001.csv}
    % ---- PTB / char-LSTM
    \nextgroupplot[title={PTB / char-LSTM, $\varepsilon=0.1$}]
    \zooplots{time}{data/curve_lstm_ptb_lr0.1.csv}
    \nextgroupplot[title={PTB / char-LSTM, $\varepsilon=0.01$}]
    \zooplots{time}{data/curve_lstm_ptb_lr0.01.csv}
    \nextgroupplot[title={PTB / char-LSTM, $\varepsilon=0.001$}]
    \zooplots{time}{data/curve_lstm_ptb_lr0.001.csv}
    % ---- PTB / char-Transformer
    \nextgroupplot[title={PTB / char-Transformer, $\varepsilon=0.1$}]
    \zooplots{time}{data/curve_transformer_ptb_lr0.1.csv}
    \nextgroupplot[title={PTB / char-Transformer, $\varepsilon=0.01$}]
    \zooplots{time}{data/curve_transformer_ptb_lr0.01.csv}
    \nextgroupplot[title={PTB / char-Transformer, $\varepsilon=0.001$}]
    \zooplots{time}{data/curve_transformer_ptb_lr0.001.csv}
  \end{groupplot}
\end{tikzpicture}

  \\[0.5ex]
  \ref{zoolegend}
  \caption{Test loss versus training time for every run. Rows: tasks; columns:
  $\varepsilon\in\{0.1,0.01,0.001\}$. Diverged runs end at their last finite
  evaluation.}
  \label{fig:curves-all}
\end{figure}

% ====================================================================================
\clearpage
\section*{NeurIPS Paper Checklist}

\begin{enumerate}

\item {\bf Claims}
    \item[] Question: Do the main claims made in the abstract and introduction accurately reflect the paper's contributions and scope?
    \item[] Answer: \answerYes{}
    \item[] Justification: The abstract and Section~\ref{sec:intro} claim an empirical comparison of eight training rules on four small architectures under one protocol; this is exactly what Sections~\ref{sec:methods}--\ref{sec:results} deliver, and the scope (single seed, one hour, $\approx$40k parameters) is stated in Section~\ref{sec:discussion}.
    \item[] Guidelines:
    \begin{itemize}
        \item The answer NA means that the abstract and introduction do not include the claims made in the paper.
        \item The abstract and/or introduction should clearly state the claims made, including the contributions made in the paper and important assumptions and limitations. A No or NA answer to this question will not be perceived well by the reviewers.
        \item The claims made should match theoretical and experimental results, and reflect how much the results can be expected to generalize to other settings.
        \item It is fine to include aspirational goals as motivation as long as it is clear that these goals are not attained by the paper.
    \end{itemize}

\item {\bf Limitations}
    \item[] Question: Does the paper discuss the limitations of the work performed by the authors?
    \item[] Answer: \answerYes{}
    \item[] Justification: Section~\ref{sec:discussion} discusses the single seed, the one-hour budget, the 40k-parameter scale, the tying of $\varepsilon$ to the learning rate, the dependence of wall-clock results on implementation efficiency, the status of \FF{} and \BP{} as non-probe references, and the character-level PTB choice.
    \item[] Guidelines:
    \begin{itemize}
        \item The answer NA means that the paper has no limitation while the answer No means that the paper has limitations, but those are not discussed in the paper.
        \item The authors are encouraged to create a separate "Limitations" section in their paper.
        \item The paper should point out any strong assumptions and how robust the results are to violations of these assumptions (e.g., independence assumptions, noiseless settings, model well-specification, asymptotic approximations only holding locally). The authors should reflect on how these assumptions might be violated in practice and what the implications would be.
        \item The authors should reflect on the scope of the claims made, e.g., if the approach was only tested on a few datasets or with a few runs. In general, empirical results often depend on implicit assumptions, which should be articulated.
        \item The authors should reflect on the factors that influence the performance of the approach. For example, a facial recognition algorithm may perform poorly when image resolution is low or images are taken in low lighting. Or a speech-to-text system might not be used reliably to provide closed captions for online lectures because it fails to handle technical jargon.
        \item The authors should discuss the computational efficiency of the proposed algorithms and how they scale with dataset size.
        \item If applicable, the authors should discuss possible limitations of their approach to address problems of privacy and fairness.
        \item While the authors might fear that complete honesty about limitations might be used by reviewers as grounds for rejection, a worse outcome might be that reviewers discover limitations that aren't acknowledged in the paper. The authors should use their best judgment and recognize that individual actions in favor of transparency play an important role in developing norms that preserve the integrity of the community. Reviewers will be specifically instructed to not penalize honesty concerning limitations.
    \end{itemize}

\item {\bf Theory assumptions and proofs}
    \item[] Question: For each theoretical result, does the paper provide the full set of assumptions and a complete (and correct) proof?
    \item[] Answer: \answerNA{}
    \item[] Justification: The paper is purely empirical and contains no theoretical results; the estimator equations in Section~\ref{sec:optimizers} restate published rules with citations.
    \item[] Guidelines:
    \begin{itemize}
        \item The answer NA means that the paper does not include theoretical results.
        \item All the theorems, formulas, and proofs in the paper should be numbered and cross-referenced.
        \item All assumptions should be clearly stated or referenced in the statement of any theorems.
        \item The proofs can either appear in the main paper or the supplemental material, but if they appear in the supplemental material, the authors are encouraged to provide a short proof sketch to provide intuition.
        \item Inversely, any informal proof provided in the core of the paper should be complemented by formal proofs provided in appendix or supplemental material.
        \item Theorems and Lemmas that the proof relies upon should be properly referenced.
    \end{itemize}

    \item {\bf Experimental result reproducibility}
    \item[] Question: Does the paper fully disclose all the information needed to reproduce the main experimental results of the paper to the extent that it affects the main claims and/or conclusions of the paper (regardless of whether the code and data are provided or not)?
    \item[] Answer: \answerYes{}
    \item[] Justification: Sections~\ref{sec:methods} and~\ref{sec:setup} and Appendix~\ref{app:impl} specify the architectures, data splits, estimators, budget, optimizer, sweep, selection rule and evaluation protocol; the code and all run logs are public.
    \item[] Guidelines:
    \begin{itemize}
        \item The answer NA means that the paper does not include experiments.
        \item If the paper includes experiments, a No answer to this question will not be perceived well by the reviewers: Making the paper reproducible is important, regardless of whether the code and data are provided or not.
        \item If the contribution is a dataset and/or model, the authors should describe the steps taken to make their results reproducible or verifiable.
        \item Depending on the contribution, reproducibility can be accomplished in various ways. For example, if the contribution is a novel architecture, describing the architecture fully might suffice, or if the contribution is a specific model and empirical evaluation, it may be necessary to either make it possible for others to replicate the model with the same dataset, or provide access to the model. In general. releasing code and data is often one good way to accomplish this, but reproducibility can also be provided via detailed instructions for how to replicate the results, access to a hosted model (e.g., in the case of a large language model), releasing of a model checkpoint, or other means that are appropriate to the research performed.
        \item While NeurIPS does not require releasing code, the conference does require all submissions to provide some reasonable avenue for reproducibility, which may depend on the nature of the contribution. For example
        \begin{enumerate}
            \item If the contribution is primarily a new algorithm, the paper should make it clear how to reproduce that algorithm.
            \item If the contribution is primarily a new model architecture, the paper should describe the architecture clearly and fully.
            \item If the contribution is a new model (e.g., a large language model), then there should either be a way to access this model for reproducing the results or a way to reproduce the model (e.g., with an open-source dataset or instructions for how to construct the dataset).
            \item We recognize that reproducibility may be tricky in some cases, in which case authors are welcome to describe the particular way they provide for reproducibility. In the case of closed-source models, it may be that access to the model is limited in some way (e.g., to registered users), but it should be possible for other researchers to have some path to reproducing or verifying the results.
        \end{enumerate}
    \end{itemize}


\item {\bf Open access to data and code}
    \item[] Question: Does the paper provide open access to the data and code, with sufficient instructions to faithfully reproduce the main experimental results, as described in supplemental material?
    \item[] Answer: \answerYes{}
    \item[] Justification: The complete code, the run configurations, the per-run JSON results and the scripts that generate every table and figure are released at \url{https://github.com/Fchaubard/zero-order-ablations}; the datasets are public and downloaded by the code.
    \item[] Guidelines:
    \begin{itemize}
        \item The answer NA means that paper does not include experiments requiring code.
        \item Please see the NeurIPS code and data submission guidelines (\url{https://nips.cc/public/guides/CodeSubmissionPolicy}) for more details.
        \item While we encourage the release of code and data, we understand that this might not be possible, so “No” is an acceptable answer. Papers cannot be rejected simply for not including code, unless this is central to the contribution (e.g., for a new open-source benchmark).
        \item The instructions should contain the exact command and environment needed to run to reproduce the results. See the NeurIPS code and data submission guidelines (\url{https://nips.cc/public/guides/CodeSubmissionPolicy}) for more details.
        \item The authors should provide instructions on data access and preparation, including how to access the raw data, preprocessed data, intermediate data, and generated data, etc.
        \item The authors should provide scripts to reproduce all experimental results for the new proposed method and baselines. If only a subset of experiments are reproducible, they should state which ones are omitted from the script and why.
        \item At submission time, to preserve anonymity, the authors should release anonymized versions (if applicable).
        \item Providing as much information as possible in supplemental material (appended to the paper) is recommended, but including URLs to data and code is permitted.
    \end{itemize}


\item {\bf Experimental setting/details}
    \item[] Question: Does the paper specify all the training and test details (e.g., data splits, hyperparameters, how they were chosen, type of optimizer, etc.) necessary to understand the results?
    \item[] Answer: \answerYes{}
    \item[] Justification: Section~\ref{sec:setup} gives the data splits, batch sizes, initialisation, optimizer, the $\varepsilon$ sweep and the validation-based selection rule; Appendix~\ref{app:impl} (Table~\ref{tab:hparams}) lists every hyper-parameter.
    \item[] Guidelines:
    \begin{itemize}
        \item The answer NA means that the paper does not include experiments.
        \item The experimental setting should be presented in the core of the paper to a level of detail that is necessary to appreciate the results and make sense of them.
        \item The full details can be provided either with the code, in appendix, or as supplemental material.
    \end{itemize}

\item {\bf Experiment statistical significance}
    \item[] Question: Does the paper report error bars suitably and correctly defined or other appropriate information about the statistical significance of the experiments?
    \item[] Answer: \answerNo{}
    \item[] Justification: Every configuration is run with a single seed because the full grid already requires 96 GPU-hours; Section~\ref{sec:discussion} states this explicitly and warns that small differences should not be read as rankings.
    \item[] Guidelines:
    \begin{itemize}
        \item The answer NA means that the paper does not include experiments.
        \item The authors should answer "Yes" if the results are accompanied by error bars, confidence intervals, or statistical significance tests, at least for the experiments that support the main claims of the paper.
        \item The factors of variability that the error bars are capturing should be clearly stated (for example, train/test split, initialization, random drawing of some parameter, or overall run with given experimental conditions).
        \item The method for calculating the error bars should be explained (closed form formula, call to a library function, bootstrap, etc.)
        \item The assumptions made should be given (e.g., Normally distributed errors).
        \item It should be clear whether the error bar is the standard deviation or the standard error of the mean.
        \item It is OK to report 1-sigma error bars, but one should state it. The authors should preferably report a 2-sigma error bar than state that they have a 96\% CI, if the hypothesis of Normality of errors is not verified.
        \item For asymmetric distributions, the authors should be careful not to show in tables or figures symmetric error bars that would yield results that are out of range (e.g. negative error rates).
        \item If error bars are reported in tables or plots, The authors should explain in the text how they were calculated and reference the corresponding figures or tables in the text.
    \end{itemize}

\item {\bf Experiments compute resources}
    \item[] Question: For each experiment, does the paper provide sufficient information on the computer resources (type of compute workers, memory, time of execution) needed to reproduce the experiments?
    \item[] Answer: \answerYes{}
    \item[] Justification: Section~\ref{sec:setup} states that every run uses one RTX~5090 for one hour of training (plus evaluation), 96 runs in total executed 8 at a time (10.5 hours of wall-clock on the 8-GPU node, about 100 GPU-hours including evaluation); the 32 two-minute smoke runs and the four supplementary 10-minute runs add about 2 GPU-hours, and the gradient-alignment measurements ran on a laptop CPU.
    \item[] Guidelines:
    \begin{itemize}
        \item The answer NA means that the paper does not include experiments.
        \item The paper should indicate the type of compute workers CPU or GPU, internal cluster, or cloud provider, including relevant memory and storage.
        \item The paper should provide the amount of compute required for each of the individual experimental runs as well as estimate the total compute.
        \item The paper should disclose whether the full research project required more compute than the experiments reported in the paper (e.g., preliminary or failed experiments that didn't make it into the paper).
    \end{itemize}

\item {\bf Code of ethics}
    \item[] Question: Does the research conducted in the paper conform, in every respect, with the NeurIPS Code of Ethics \url{https://neurips.cc/public/EthicsGuidelines}?
    \item[] Answer: \answerYes{}
    \item[] Justification: The study trains small models on three standard public benchmarks; it involves no human subjects, no personal data and no deployment.
    \item[] Guidelines:
    \begin{itemize}
        \item The answer NA means that the authors have not reviewed the NeurIPS Code of Ethics.
        \item If the authors answer No, they should explain the special circumstances that require a deviation from the Code of Ethics.
        \item The authors should make sure to preserve anonymity (e.g., if there is a special consideration due to laws or regulations in their jurisdiction).
    \end{itemize}


\item {\bf Broader impacts}
    \item[] Question: Does the paper discuss both potential positive societal impacts and negative societal impacts of the work performed?
    \item[] Answer: \answerNA{}
    \item[] Justification: This is foundational optimizer-benchmarking research on $\approx$40k-parameter models with no direct application or deployment path; its potential positive impact (more energy-efficient or hardware-friendly training) is discussed in Section~\ref{sec:intro} and we see no specific negative impact beyond those of generic training algorithms.
    \item[] Guidelines:
    \begin{itemize}
        \item The answer NA means that there is no societal impact of the work performed.
        \item If the authors answer NA or No, they should explain why their work has no societal impact or why the paper does not address societal impact.
        \item Examples of negative societal impacts include potential malicious or unintended uses (e.g., disinformation, generating fake profiles, surveillance), fairness considerations (e.g., deployment of technologies that could make decisions that unfairly impact specific groups), privacy considerations, and security considerations.
        \item The conference expects that many papers will be foundational research and not tied to particular applications, let alone deployments. However, if there is a direct path to any negative applications, the authors should point it out. For example, it is legitimate to point out that an improvement in the quality of generative models could be used to generate deepfakes for disinformation. On the other hand, it is not needed to point out that a generic algorithm for optimizing neural networks could enable people to train models that generate Deepfakes faster.
        \item The authors should consider possible harms that could arise when the technology is being used as intended and functioning correctly, harms that could arise when the technology is being used as intended but gives incorrect results, and harms following from (intentional or unintentional) misuse of the technology.
        \item If there are negative societal impacts, the authors could also discuss possible mitigation strategies (e.g., gated release of models, providing defenses in addition to attacks, mechanisms for monitoring misuse, mechanisms to monitor how a system learns from feedback over time, improving the efficiency and accessibility of ML).
    \end{itemize}

\item {\bf Safeguards}
    \item[] Question: Does the paper describe safeguards that have been put in place for responsible release of data or models that have a high risk for misuse (e.g., pretrained language models, image generators, or scraped datasets)?
    \item[] Answer: \answerNA{}
    \item[] Justification: No models with misuse potential are released; the released artefacts are training code and logs of $\approx$40k-parameter models on MNIST, CIFAR-10 and character-level PTB.
    \item[] Guidelines:
    \begin{itemize}
        \item The answer NA means that the paper poses no such risks.
        \item Released models that have a high risk for misuse or dual-use should be released with necessary safeguards to allow for controlled use of the model, for example by requiring that users adhere to usage guidelines or restrictions to access the model or implementing safety filters.
        \item Datasets that have been scraped from the Internet could pose safety risks. The authors should describe how they avoided releasing unsafe images.
        \item We recognize that providing effective safeguards is challenging, and many papers do not require this, but we encourage authors to take this into account and make a best faith effort.
    \end{itemize}

\item {\bf Licenses for existing assets}
    \item[] Question: Are the creators or original owners of assets (e.g., code, data, models), used in the paper, properly credited and are the license and terms of use explicitly mentioned and properly respected?
    \item[] Answer: \answerYes{}
    \item[] Justification: MNIST, CIFAR-10 and Penn Treebank are cited (Section~\ref{sec:setup}) and their terms of use are stated in Appendix~\ref{app:impl}; PyTorch is used under its BSD licence.
    \item[] Guidelines:
    \begin{itemize}
        \item The answer NA means that the paper does not use existing assets.
        \item The authors should cite the original paper that produced the code package or dataset.
        \item The authors should state which version of the asset is used and, if possible, include a URL.
        \item The name of the license (e.g., CC-BY 4.0) should be included for each asset.
        \item For scraped data from a particular source (e.g., website), the copyright and terms of service of that source should be provided.
        \item If assets are released, the license, copyright information, and terms of use in the package should be provided. For popular datasets, \url{paperswithcode.com/datasets} has curated licenses for some datasets. Their licensing guide can help determine the license of a dataset.
        \item For existing datasets that are re-packaged, both the original license and the license of the derived asset (if it has changed) should be provided.
        \item If this information is not available online, the authors are encouraged to reach out to the asset's creators.
    \end{itemize}

\item {\bf New assets}
    \item[] Question: Are new assets introduced in the paper well documented and is the documentation provided alongside the assets?
    \item[] Answer: \answerYes{}
    \item[] Justification: The released code repository (MIT licence) documents the protocol (\texttt{PROTOCOL.md}), the command to launch every run and the scripts that regenerate every table and figure; the 96 result files and the supplementary runs are included.
    \item[] Guidelines:
    \begin{itemize}
        \item The answer NA means that the paper does not release new assets.
        \item Researchers should communicate the details of the dataset/code/model as part of their submissions via structured templates. This includes details about training, license, limitations, etc.
        \item The paper should discuss whether and how consent was obtained from people whose asset is used.
        \item At submission time, remember to anonymize your assets (if applicable). You can either create an anonymized URL or include an anonymized zip file.
    \end{itemize}

\item {\bf Crowdsourcing and research with human subjects}
    \item[] Question: For crowdsourcing experiments and research with human subjects, does the paper include the full text of instructions given to participants and screenshots, if applicable, as well as details about compensation (if any)?
    \item[] Answer: \answerNA{}
    \item[] Justification: The paper involves neither crowdsourcing nor human subjects.
    \item[] Guidelines:
    \begin{itemize}
        \item The answer NA means that the paper does not involve crowdsourcing nor research with human subjects.
        \item Including this information in the supplemental material is fine, but if the main contribution of the paper involves human subjects, then as much detail as possible should be included in the main paper.
        \item According to the NeurIPS Code of Ethics, workers involved in data collection, curation, or other labor should be paid at least the minimum wage in the country of the data collector.
    \end{itemize}

\item {\bf Institutional review board (IRB) approvals or equivalent for research with human subjects}
    \item[] Question: Does the paper describe potential risks incurred by study participants, whether such risks were disclosed to the subjects, and whether Institutional Review Board (IRB) approvals (or an equivalent approval/review based on the requirements of your country or institution) were obtained?
    \item[] Answer: \answerNA{}
    \item[] Justification: The paper involves no human subjects, so no IRB review was required.
    \item[] Guidelines:
    \begin{itemize}
        \item The answer NA means that the paper does not involve crowdsourcing nor research with human subjects.
        \item Depending on the country in which research is conducted, IRB approval (or equivalent) may be required for any human subjects research. If you obtained IRB approval, you should clearly state this in the paper.
        \item We recognize that the procedures for this may vary significantly between institutions and locations, and we expect authors to adhere to the NeurIPS Code of Ethics and the guidelines for their institution.
        \item For initial submissions, do not include any information that would break anonymity (if applicable), such as the institution conducting the review.
    \end{itemize}

\item {\bf Declaration of LLM usage}
    \item[] Question: Does the paper describe the usage of LLMs if it is an important, original, or non-standard component of the core methods in this research? Note that if the LLM is used only for writing, editing, or formatting purposes and does not impact the core methodology, scientific rigorousness, or originality of the research, declaration is not required.
    \item[] Answer: \answerNA{}
    \item[] Justification: LLM-based assistants were used only for writing and coding assistance; no LLM is part of the methods, the experiments or the analysis.
    \item[] Guidelines:
    \begin{itemize}
        \item The answer NA means that the core method development in this research does not involve LLMs as any important, original, or non-standard components.
        \item Please refer to our LLM policy (\url{https://neurips.cc/Conferences/2025/LLM}) for what should or should not be described.
    \end{itemize}

\end{enumerate}

\end{document}
